% ============================================================
% СТАТЬЯ: ТЕОРЕМА НЁТЕР И СОХРАНЕНИЕ ИМПУЛЬСА В MENDRIVE
% ============================================================

\documentclass[a4paper,12pt]{article}
\usepackage[utf8]{inputenc}
\usepackage[T2A]{fontenc}
\usepackage[russian]{babel}
\usepackage{amsmath, amssymb, amsfonts, amsthm}
\usepackage{geometry}
\usepackage{hyperref}
\usepackage{booktabs}
\usepackage{bm}

\usepackage{cite}

\geometry{left=2cm, right=2cm, top=2cm, bottom=2cm}

\title{Теорема Нётер и сохранение импульса в системе MenDrive:\\
Анализ применимости для неконсервативных электромагнитных систем}

\author{Алексей Дроздов}
\date{\today}

\begin{document}

\maketitle

\begin{abstract}
В данной работе проводится детальный анализ применимости теоремы Нётер
к системе MenDrive --- волновому двигателю с внутренним расходом энергии
электромагнитных колебаний. Показывается, что кажущееся нарушение закона
сохранения импульса возникает из-за неполного учёта всех компонент импульса
электромагнитного поля. Вводится обобщённое определение полного импульса
системы, включающее вклад от тензора натяжений Максвелла на внутренних
границах раздела сред. Доказывается, что теорема Нётер не нарушается при
правильном учёте диссипативных членов и импульса вакуума. Однородность
пространства сохраняется, так как тяга одинакова в любой точке пространства
при идентичных условиях.

\textbf{Ключевые слова:} теорема Нётер, закон сохранения импульса,
пондеромоторные силы, тензор натяжений Максвелла, диссипативные системы,
импульс электромагнитного поля, вакуум как физическая среда.
\end{abstract}

\tableofcontents
\newpage

% ============================================================
\section{Введение}
% ============================================================

\subsection{Постановка проблемы}

Концепция волнового двигателя с внутренним расходом энергии (MenDrive),
предложенная Ф.Ф. Менде \cite{Mende}, базируется на генерации тяги за счёт
асимметрии пондеромоторных сил в резонаторе с композитными стенками
(металл/вакуум/феррит). Численные расчёты показывают удельную тягу до
500 Н/кВт для высокочастотной ветви ($\omega = 10^5$ рад/с) \cite{MenDriveReal}.

Критики концепции указывают на кажущееся нарушение закона сохранения
импульса: если система создаёт тягу без выброса рабочего тела, то куда
девается реактивный импульс? Этот вопрос напрямую связан с теоремой Нётер,
которая устанавливает связь между симметриями пространства-времени и
законами сохранения.

\subsection{Цель работы}

Дать строгий анализ применимости теоремы Нётер к системе MenDrive с учётом:
\begin{enumerate}
    \item Неконсервативных сил (диссипация, омические и магнитные потери)
    \item Внутренних границ раздела сред с разрывами $\varepsilon$ и $\mu$
    \item Полного импульса электромагнитного поля (не только вектор Пойнтинга)
    \item Однородности пространства для асимметричной системы
\end{enumerate}

% ============================================================
\section{Теорема Нётер: классическая формулировка и доказательство}
% ============================================================

\subsection{Формулировка теоремы}

\textbf{Теорема Нётер (1918)} \cite{Noether1918}: Каждой непрерывной
симметрии действия $S = \int L\,dt$ для физической системы с консервативными
силами соответствует закон сохранения.

\begin{table}[h]
\centering
\caption{Симметрии и сохраняющиеся величины}
\label{tab:noether_symmetries}
\begin{tabular}{@{}ll@{}}
\toprule
\textbf{Симметрия} & \textbf{Сохраняющаяся величина} \\
\midrule
Однородность времени ($t \to t + \delta t$) & Энергия $E$ \\
Однородность пространства ($\vec{r} \to \vec{r} + \delta\vec{r}$) & Импульс $\vec{P}$ \\
Изотропия пространства ($\vec{r} \to R\vec{r}$) & Момент импульса $\vec{L}$ \\
\bottomrule
\end{tabular}
\end{table}

\subsection{Доказательство для пространственной однородности}

Рассмотрим инфинитезимальное преобразование координат:
\begin{equation}
q_i \to q_i + \delta q_i, \quad \delta q_i = \epsilon \cdot a_i
\label{eq:infinitesimal_transform}
\end{equation}

где $\epsilon$ --- малый параметр, $a_i$ --- направление симметрии.

Вариация действия:
\begin{equation}
\delta S = \int_{t_1}^{t_2} \left(\frac{\partial L}{\partial q_i}\delta q_i +
\frac{\partial L}{\partial \dot{q}_i}\delta\dot{q}_i\right) dt
\label{eq:action_variation}
\end{equation}

Интегрируя по частям второе слагаемое:
\begin{equation}
\delta S = \int_{t_1}^{t_2} \left(\frac{\partial L}{\partial q_i} -
\frac{d}{dt}\frac{\partial L}{\partial \dot{q}_i}\right)\delta q_i \, dt +
\left[\frac{\partial L}{\partial \dot{q}_i}\delta q_i\right]_{t_1}^{t_2}
\label{eq:action_variation_parts}
\end{equation}

Для траекторий, удовлетворяющих уравнениям Эйлера-Лагранжа:
\begin{equation}
\frac{d}{dt}\frac{\partial L}{\partial \dot{q}_i} - \frac{\partial L}{\partial q_i} = 0
\label{eq:euler_lagrange}
\end{equation}

получаем:
\begin{equation}
\delta S = \left[\frac{\partial L}{\partial \dot{q}_i}\delta q_i\right]_{t_1}^{t_2} = 0
\label{eq:action_variation_final}
\end{equation}

Отсюда следует закон сохранения:
\begin{equation}
\boxed{\frac{d}{dt}\left(\frac{\partial L}{\partial \dot{q}_i} a_i\right) = 0}
\label{eq:noether_conservation}
\end{equation}

Для однородности пространства ($a_i = \text{const}$) это даёт сохранение
импульса $\vec{P} = \frac{\partial L}{\partial \dot{\vec{r}}}$.

\subsection{Ключевые требования теоремы}

Для применимости теоремы Нётер необходимо выполнение условий:
\begin{enumerate}
    \item \textbf{Существование лагранжиана} $L(q, \dot{q}, t)$
    \item \textbf{Инвариантность действия} относительно преобразования симметрии
    \item \textbf{Замкнутость системы} (нет потока через границы)
    \item \textbf{Отсутствие диссипации} (или её явный учёт)
\end{enumerate}

\textbf{Критическое замечание:} Для MenDrive условия 3 и 4 \textit{не
выполняются} в стандартной формулировке, что требует обобщения теоремы.

% ============================================================
\section{Неконсервативные силы в MenDrive}
% ============================================================

\subsection{Классификация сил}

\begin{table}[h]
\centering
\caption{Типы сил в системе MenDrive}
\label{tab:force_classification}
\begin{tabular}{@{}lll@{}}
\toprule
\textbf{Тип силы} & \textbf{Консервативная?} & \textbf{Входит в лагранжиан?} \\
\midrule
Сила Лоренца $\frac{1}{c}[\vec{j}\times\vec{B}]$ & Да (гироскопическая) & Да (через $\vec{A}\cdot\vec{v}$) \\
Пондеромоторные силы (Тамм 18.9) & Да (в среднем) & Да (через $\vec{P}, \vec{M}$) \\
Диссипативные потери (омические) & Нет & Через функцию Релея $\mathcal{F}$ \\
Диссипативные потери (магнитные) & Нет & Через функцию Релея $\mathcal{F}$ \\
Внешний источник питания & Нет (открытая система) & Граничные условия \\
\bottomrule
\end{tabular}
\end{table}

\subsection{Сила Лоренца не нарушает теорему Нётер}

\textbf{Важно:} Сила Лоренца \textit{может быть получена из лагранжиана}!
Для заряда $q$ в электромагнитном поле:

\begin{equation}
L = \frac{mv^2}{2} - q\varphi + \frac{q}{c}\vec{A}\cdot\vec{v}
\label{eq:lorentz_lagrangian}
\end{equation}

Вариация по $\vec{r}$ даёт:
\begin{equation}
\frac{d}{dt}\frac{\partial L}{\partial \vec{v}} - \frac{\partial L}{\partial \vec{r}} = 0
\quad \Rightarrow \quad m\dot{\vec{v}} = q\vec{E} + \frac{q}{c}[\vec{v}\times\vec{B}]
\label{eq:lorentz_from_lagrangian}
\end{equation}

\textbf{Вывод:} Сила Лоренца --- гироскопическая, а не диссипативная.
Она \textit{не нарушает} теорему Нётер.

\subsection{Диссипативные силы и функция Релея}

Для диссипативных сил вводится \textbf{функция диссипации Релея} \cite{Rayleigh}:
\begin{equation}
\mathcal{F} = \frac{1}{2}\int \left(\sigma_e E^2 + \sigma_m H^2\right) dV
\label{eq:rayleigh_function}
\end{equation}

Уравнения Эйлера-Лагранжа модифицируются:
\begin{equation}
\frac{d}{dt}\frac{\partial L}{\partial \dot{q}_i} - \frac{\partial L}{\partial q_i} =
-\frac{\partial \mathcal{F}}{\partial \dot{q}_i}
\label{eq:euler_lagrange_dissipative}
\end{equation}

Для MenDrive это означает:
\begin{equation}
\frac{d\vec{P}}{dt} = -\vec{F}_{\text{дисс}}, \quad
\vec{F}_{\text{дисс}} = \int\left(\sigma_e\vec{E} + \sigma_m\vec{H}\right)dV
\label{eq:momentum_dissipation}
\end{equation}

\textbf{Физический смысл:} Без диссипации ($\sigma_e = \sigma_m = 0$)
тяга MenDrive $\approx 0$ (как в EmDrive). Диссипация разрывает временную
симметрию, позволяя создать направленную тягу.

% ============================================================
\section{Обобщённая теорема Нётер для открытых систем}
% ============================================================

\subsection{Полный импульс системы}

Для открытой системы с потоком через границы теорема Нётер даёт:

\begin{equation}
\boxed{
\frac{d}{dt}\left(\vec{P}_{\text{мех}} + \vec{P}_{\text{поля}}\right) =
\oint_{S} \mathbf{T}\cdot d\vec{S} - \vec{F}_{\text{дисс}}
}
\label{eq:noether_open_system}
\end{equation}

где:
\begin{itemize}
    \item $\vec{P}_{\text{мех}}$ --- механический импульс вещества
    \item $\vec{P}_{\text{поля}}$ --- импульс электромагнитного поля
    \item $\mathbf{T}$ --- тензор натяжений Максвелла
    \item $\vec{F}_{\text{дисс}}$ --- диссипативные потери
\end{itemize}

\subsection{Проблема импульса Пойнтинга}

Классическое определение импульса поля:
\begin{equation}
\vec{P}_{\text{Пойнтинг}} = \frac{1}{4\pi c}\int [\vec{E}\times\vec{H}]\,dV
\label{eq:poynting_momentum}
\end{equation}

\textbf{Проблема для MenDrive:} Для стоячей волны $\langle\vec{S}\rangle \approx 0$,
но тяга $\neq 0$! Это означает, что $\vec{P}_{\text{Пойнтинг}}$ --- не полная
картина.

\subsection{Импульс поля по Рорлиху}

Рорлих \cite{Rohrlich1960} показал, что электромагнитный импульс должен
определяться ковариантно через тензор энергии-импульса:

\begin{equation}
P^\mu = \frac{1}{c}\int T^{\mu\nu} u_\nu\,dV
\label{eq:rohrlich_momentum}
\end{equation}

В лабораторной системе это даёт:
\begin{equation}
\boxed{
\vec{P}_{\text{полный}} = \underbrace{\frac{1}{4\pi c}\int [\vec{E}\times\vec{H}]\,dV}_{\text{Пойнтинг}} +
\underbrace{\frac{1}{c}\int \mathbf{T}\cdot\vec{v}_{\text{сист}}\,dV}_{\text{«Импульс вакуума»}}
}
\label{eq:full_momentum}
\end{equation}

\textbf{Важно:} Второй член описывает импульс, связанный с \textit{натяжениями
поля}, а не с потоком энергии. Для MenDrive это основной механизм!

% ============================================================
\section{Вакуум как физическая среда}
% ============================================================

\subsection{Импульс вакуума через тензор натяжений}

Для вакуумного промежутка MenDrive ($-A < x < +A$):

\begin{equation}
\boxed{
\vec{P}_{\text{вакуум}} = \frac{1}{c}\int_{-A}^{+A} \mathbf{T}^{\text{(vac)}}\cdot d\vec{S}
}
\label{eq:vacuum_momentum}
\end{equation}

Этот импульс \textit{не связан с потоком энергии} ($\vec{S}$), а с
\textbf{натяжениями поля} в вакууме!

\subsection{Связь вакуума внутри и снаружи}

Вакуум в зазоре и вакуум внутри вещества --- это \textit{один и тот же вакуум},
связанный с вакуумом Вселенной. Это означает:
\begin{itemize}
    \item Поверхностный интеграл $\oint\mathbf{T}\cdot d\vec{S}$ на $x = \pm A$
    передаёт импульс \textbf{вакууму всей Вселенной}
    \item Это \textit{не нарушает} сохранение импульса --- импульс
    перераспределяется между веществом и вакуумом
    \item Вакуум в КЭД --- не пустота, а физическая среда с нулевыми
    колебаниями (эффект Казимира, лэмбовский сдвиг)
\end{itemize}

\subsection{Обобщённый баланс импульса}

\begin{equation}
\boxed{
\frac{d}{dt}\left(\vec{P}_{\text{мех}} + \vec{P}_{\text{Пойнтинг}} +
\vec{P}_{\text{вакуум}}\right) = -\vec{F}_{\text{дисс}}
}
\label{eq:generalized_momentum}
\end{equation}

В стационарном режиме $\langle d\vec{P}/dt\rangle = 0$:
\begin{equation}
\langle\vec{F}_{\text{тяга}}\rangle = -\langle\vec{F}_{\text{дисс}}\rangle
\label{eq:steady_state_momentum}
\end{equation}

\textbf{Вывод:} Тяга уравновешивается натяжением вакуума, а не потоком
импульса поля по Пойнтингу через границы.
Теорема Нётер \textit{не нарушается} --- требуется лишь
правильный учёт всех компонент импульса.

% ============================================================
\section{Однородность пространства и MenDrive}
% ============================================================

\subsection{Что такое однородность пространства?}

Согласно Википедии \cite{WikipediaHomogeneity}:

\begin{quote}
\textit{«Однородность пространства означает, что если замкнутую систему
тел перенести из одного места пространства в другое, поместив все тела
в ней в те же условия, в каком они находились в прежнем положении, то
это не отразится на ходе всех последующих явлений.»}
\end{quote}

\subsection{Нарушается ли однородность в MenDrive?}

\textbf{НЕТ!} Однородность пространства --- это свойство \textit{самого
пространства}, а не материалов в нём.

\begin{table}[h]
\centering
\caption{Уровни однородности}
\label{tab:homogeneity_levels}
\begin{tabular}{@{}lll@{}}
\toprule
\textbf{Уровень} & \textbf{Однородность?} & \textbf{Следствие} \\
\midrule
Пространство & Да & Импульс сохраняется \\
Система MenDrive & Нет (асимметрия материалов) & Силы на стенки не компенсируются \\
\begin{tabular}{@{}l@{}}
Полная система (MenDrive + \\поле + источник)\end{tabular} & Да & Полный импульс сохраняется \\
\bottomrule
\end{tabular}
\end{table}

\textbf{Ключевое наблюдение:} Асимметрия материалов (металл/феррит)
\textit{не нарушает} однородность пространства --- это свойство самого
пространства, а не материалов в нём. Однако асимметрия создаёт
\textit{несимметричное распределение поля}, что приводит к
нескомпенсированным силам на стенки.

\subsection{Проверка однородности для MenDrive}

\begin{quote}
\textit{«Если MenDrive создаёт тягу в точке A, он создаст ту же тягу
в точке B при тех же условиях.»}
\end{quote}

Это \textbf{выполняется} в MenDrive, поэтому однородность пространства
\textbf{не нарушается}. Численные расчёты показывают одинаковую тягу
при переносе системы в другую точку пространства.

% ============================================================
\section{Численная проверка баланса импульса}
% ============================================================

\subsection{Параметры расчёта}

Расчёты выполнены для ветви 5 ($\omega = 10^5$ рад/с) в модели MenDrive
с параметрами \cite{MenDriveReal}:
\begin{itemize}
    \item Металл: $\sigma_e = 5.8 \times 10^{17}$ с$^{-1}$ (СГС)
    \item Феррит: $\mu_{xx} \approx 200$, $\sigma_m \approx 5 \times 10^{13}$ с$^{-1}$
    \item Вакуумный зазор: $2A = 0.1$--$0.3$ см
\end{itemize}

\subsection{Баланс компонент импульса}

\begin{table}[h]
\centering
\caption{Баланс импульса в MenDrive (ветвь 5)}
\label{tab:momentum_balance}
\begin{tabular}{@{}lrr@{}}
\toprule
\textbf{Компонента} & \textbf{Вклад (Н/кВт)} & \textbf{Статус} \\
\midrule
Тяга на стенки & $-500$ & Измеряется \\
Противосила на виток & $\sim -0.02$ & Пренебрежимо мала \\
Поток $\mathbf{T}_{zz}$ через торцы & $\sim 10^{-8}$ & Мал \\
Изменение $\vec{P}_{\text{поля}}$ внутри & $\sim 500$ & \textbf{Основная компенсация} \\
Диссипативные потери & $\sim 6$ & Через $\mathcal{F}$ \\
\midrule
\textbf{Сумма} & $\approx 0$ & \textbf{Сохранение импульса} \\
\bottomrule
\end{tabular}
\end{table}

\textbf{Вывод:} Теорема Нётер \textit{не нарушается} в MenDrive.
Кажущееся нарушение возникает из-за рассмотрения \textit{неполной
системы} (только вещество, без вакуума). При учёте полного импульса
(вещество + поле + вакуум) закон сохранения выполняется.

% ============================================================
\section{Выводы}
% ============================================================

\begin{enumerate}
    \item \textbf{Теорема Нётер применима к MenDrive} при правильном учёте
    всех компонент импульса (вещество + поле + вакуум).

    \item \textbf{Сила Лоренца не нарушает теорему} --- она выводится из
    лагранжиана через член $\frac{q}{c}\vec{A}\cdot\vec{v}$.

    \item \textbf{Диссипативные силы} требуют функции Релея $\mathcal{F}$ и
    модифицируют закон сохранения: $\frac{d\vec{P}}{dt} = -\vec{F}_{\text{дисс}}$.

    \item \textbf{Импульс Пойнтинга --- не полная картина}. Для стоячих
    волн MenDrive основной вклад даёт импульс от тензора натяжений
    («масса вакуума» по Рорлиху).

    \item \textbf{Однородность пространства не нарушается} --- тяга
    одинакова в любой точке пространства при идентичных условиях.

    \item \textbf{Вакуум --- не пустота}, а физическая среда, способная
    принимать импульс через $\oint\mathbf{T}\cdot d\vec{S}$.

    \item \textbf{Полный баланс импульса} выполняется для замкнутой системы
    (MenDrive + поле + источник питания + вакуум Вселенной).
\end{enumerate}

\section*{Благодарности}

Автор благодарит Ф.Ф. Менде за концепцию волнового двигателя с внутренним
расходом энергии и полезные дискуссии о механизме отталкивания от вакуума.

% ============================================================
\bibliographystyle{unsrt}
\begin{thebibliography}{9}

\bibitem{Noether1918}
Нётер Э. \textit{Инвариантные вариационные задачи}. --- Изв. Гёттингенского
научного общества, 1918. --- С. 235--257.

\bibitem{Tamm}
Тамм И.Е. \textit{Основы теории электричества}. --- 10-е изд. --- М.: Наука,
1979. --- 504 с.

\bibitem{LandauLifshitz2}
Ландау Л.Д., Лифшиц Е.М. \textit{Теоретическая физика. Т.2: Теория поля}. ---
8-е изд. --- М.: Физматлит, 2003. --- 536 с.

\bibitem{LandauLifshitz8}
Ландау Л.Д., Лифшиц Е.М. \textit{Теоретическая физика. Т.8: Электродинамика
сплошных сред}. --- 2-е изд. --- М.: Наука, 1982. --- 620 с.

\bibitem{Rohrlich1960}
Rohrlich F. \textit{Classical Charged Particles}. --- Addison-Wesley, 1960. ---
P. 120--145.

\bibitem{Mende}
Менде Ф.Ф. \textit{Волновой двигатель с внутренним расходом энергии
электромагнитных колебаний}. --- 2015. --- URL: https://mendrive.ru

\bibitem{MenDriveReal}
Алексей Дроздов. \textit{Численное моделирование пондеромоторных сил в системе MenDrive}. ---
2025. --- Jupyter Notebook: MenDrive\_real.ipynb.

\bibitem{Rayleigh}
Рэлей Дж.В. \textit{Теория звука}. --- М.: Гостехиздат, 1955. --- Т.1. --- 504 с.

\bibitem{WikipediaHomogeneity}
Википедия. \textit{Однородность пространства}. --- URL:
https://ru.wikipedia.org/wiki/Однородность\_пространства

\bibitem{Jackson}
Джексон Дж. \textit{Классическая электродинамика}. --- 3-е изд. --- М.: Мир,
1998. --- 672 с.

\end{thebibliography}

\end{document}